\subsection{Conceptual Implications and Open Challenges}
  \label{subsec:conceptual-implications-and-open-challenges}

  Temporal ordering arises from the monotone growth of cumulative projected
  capacity $I(n)$ along the admissible cascade~\cite{Beau2026pto};
  energy quantifies residual projective capacity available for further structural
  advancement; irreversibility reflects the one-way activation of projected modes.
  Time, energy, and irreversibility are not independent axioms but complementary
  effective descriptions of the same projective structure.

  Open challenges include the quantitative reconstruction of CMB anisotropies from
  early-stage projective ordering structure, the detailed treatment of non-equilibrium
  decoherence and reprojection, the emergence of gauge symmetries from topological
  features, the explicit derivation of the full flavor hierarchy and the quantitative
  determination of the CP-violating invariant from first-principles spectral admissibility
  structure, and the long-term stability of solitonic configurations under extreme
  conditions.
  Progress will require large-scale numerical campaigns on the spectral admissibility
  pipeline, discretized lattice realizations, and targeted experimental tests.

  \paragraph{Possible Cosmological Implications of Flavor Misalignment.}

    Within the present framework, the Jarlskog invariant $\mathcal{J}$ is interpreted
    as a structural residue of quasi-commutative sectorial restrictions of the admissible
    spectral subspace of the projection fiber~$\Pi$
    (Appendix~\ref{subsec:cp-asymmetry-and-chiral-selection}).
    In the effective low-energy regime, this invariant is treated as
    a stable characteristic of the stabilized flavor manifold.

    Because the operators $\Pi_u$ and $\Pi_d$ are defined as
    restrictions of the global non-injective mapping,
    their algebraic structure is ultimately rooted in the global
    state of the relational substrate~$\chi$.
    In extreme regimes approaching global saturation---such as the
    early Universe---the admissible spectral subspace of the fiber
    could differ from its present quasi-stationary configuration.

    Whether such regimes might modify the effective structure of
    flavor misalignment remains an open question.
    No quantitative model is presently derived.
    Establishing a predictive relation between $\mathcal{J}$ and
    early-Universe projective ordering conditions constitutes a
    significant theoretical challenge for future work.
